\chapter{Build, imagem e execu\c{c}\~ao}
\label{cap:build}

\section{Toolchain e alvo}

O projeto compila para \code{x86\_64-unknown-none}, alvo oficial de Rust para
bare metal x86\_64. Esse alvo:

\begin{itemize}
  \item n\~ao possui \code{std}, \code{alloc} nem \emph{runtime};
  \item usa o modelo de c\'odigo \emph{kernel}, que assume endere\c{c}os no
        ter\c{c}o superior do espa\c{c}o de endere\c{c}amento;
  \item produz c\'odigo posicional apenas quando necess\'ario, com
        \code{relocation-model=static}.
\end{itemize}

O arquivo \code{rust-toolchain.toml} fixa o canal \code{stable} e solicita o
alvo:

\begin{lstlisting}[caption={\code{rust-toolchain.toml}.},label={lst:toolchain}]
[toolchain]
channel = "stable"
targets = ["x86_64-unknown-none"]
\end{lstlisting}

\section{Configura\c{c}\~ao do Cargo}

\begin{lstlisting}[caption={\code{.cargo/config.toml}.},label={lst:cargoconfig}]
[build]
target = "x86_64-unknown-none"

[target.x86_64-unknown-none]
rustflags = [
    "-C", "relocation-model=static",
    "-C", "linker=ld.lld",
    "-C", "link-arg=-Tlinker.ld",
    "-C", "link-arg=--gc-sections",
]
\end{lstlisting}

\begin{table}[H]
\centering
\caption{Flags de compila\c{c}\~ao e liga\c{c}\~ao.}
\label{tab:flags}
\begin{tabularx}{\linewidth}{@{}lX@{}}
\toprule
\textbf{Flag} & \textbf{Efeito} \\
\midrule
\code{relocation-model=static} & Evita \emph{PIC}/GOT; o kernel \'e ligado em
endere\c{c}o fixo. \\
\code{linker=ld.lld} & Usa o linker LLVM, compat\'ivel com o script. \\
\code{-Tlinker.ld} & Aplica o layout de mem\'oria do projeto. \\
\code{--gc-sections} & Remove se\c{c}\~oes n\~ao referenciadas. \\
\bottomrule
\end{tabularx}
\end{table}

\section{Perfis de compila\c{c}\~ao}

\begin{lstlisting}[caption={\code{Cargo.toml}.},label={lst:cargo}]
[package]
name = "vga-kernel"
version = "0.1.0"
edition = "2021"
description = "Kernel minimalista em Rust: boot BIOS (multiboot2), console VGA 80x25 e teclado PS/2"
license = "MIT OR Apache-2.0"

[dependencies]

[profile.dev]
panic = "abort"

[profile.release]
panic = "abort"
opt-level = 2
lto = true
codegen-units = 1
\end{lstlisting}

O perfil \code{release} \'e usado na ISO. \code{panic=abort} \'e obrigat\'orio
para um kernel sem \emph{unwinding}; \code{lto} e \code{codegen-units=1}
permitem otimiza\c{c}\~oes entre m\'odulos e reduzem o tamanho.

\section{Makefile}

\begin{lstlisting}[caption={Alvos principais do \code{Makefile}.},label={lst:makefile}]
TARGET  := x86_64-unknown-none
KERNEL  := target/$(TARGET)/release/vga-kernel
ISO     := vga-kernel.iso
ISODIR  := build/iso

RUSTC   ?= $(shell command -v rustup >/dev/null 2>&1 && rustup which rustc)
ifneq ($(RUSTC),)
export RUSTC
endif

CARGO   ?= $(shell command -v rustup >/dev/null 2>&1 && echo "rustup run stable cargo" || echo cargo)

all: iso

kernel:
	$(CARGO) build --release --target $(TARGET)

iso: kernel
	rm -rf $(ISODIR)
	mkdir -p $(ISODIR)/boot/grub
	cp $(KERNEL) $(ISODIR)/boot/kernel.elf
	cp boot/grub/grub.cfg $(ISODIR)/boot/grub/grub.cfg
	grub-mkrescue -o $(ISO) $(ISODIR)

run: iso
	qemu-system-x86_64 -cdrom $(ISO) -boot order=d -m 128M -serial stdio

clean:
	$(CARGO) clean
	rm -rf build $(ISO)
\end{lstlisting}

\begin{nota}[Compatibilidade com o sistema]
Em algumas distribui\c{c}\~oes, \code{/usr/bin/cargo} e \code{/usr/bin/rustc}
n\~ao s\~ao os do \code{rustup} e n\~ao t\^em o alvo
\code{x86\_64-unknown-none}. O Makefile detecta \code{rustup}, usa
\code{rustup run stable cargo} e exporta \code{RUSTC} para que o Cargo n\~ao
encontre o \code{rustc} errado pelo \code{PATH}.
\end{nota}

\subsection{Alvos}

\begin{table}[H]
\centering
\caption{Alvos do Makefile.}
\label{tab:maketargets}
\begin{tabularx}{\linewidth}{@{}lX@{}}
\toprule
\textbf{Alvo} & \textbf{A\c{c}\~ao} \\
\midrule
\code{make} / \code{make all} & Compila o kernel e gera a ISO. \\
\code{make kernel} & Compila apenas o ELF release. \\
\code{make iso} & Compila e empacota a ISO com GRUB. \\
\code{make run} & Gera a ISO e inicia o QEMU com VGA e serial no stdio. \\
\code{make clean} & Remove \code{target/}, \code{build/} e a ISO. \\
\bottomrule
\end{tabularx}
\end{table}

\section{Configura\c{c}\~ao do GRUB}

\begin{lstlisting}[caption={\code{boot/grub/grub.cfg}.},label={lst:grubcfg}]
set timeout=0
set default=0

terminal_output console
set gfxpayload=text

insmod multiboot2

menuentry "Rust VGA Kernel (x86_64)" {
    multiboot2 /boot/kernel.elf
    boot
}
\end{lstlisting}

\begin{itemize}
  \item \code{terminal\_output console} mant\'em o GRUB no modo texto VGA,
        garantindo que o kernel encontre \code{0xB8000} como esperado.
  \item \code{gfxpayload=text} impede que o GRUB selecione um modo gr\'afico
        linear.
  \item \code{insmod multiboot2} carrega o m\'odulo do protocolo.
  \item \code{multiboot2} instrui o GRUB a carregar o ELF e passar a
        assinatura e a estrutura de informa\c{c}\~oes.
\end{itemize}

\section{ISO}

O alvo \code{iso} monta uma \'arvore tempor\'aria:

\begin{lstlisting}[style=shell,caption={Conte\'udo da \'arvore da ISO.},label={lst:isotree}]
build/iso/
|-- boot/
|   |-- grub/
|   |   `-- grub.cfg
|   `-- kernel.elf
`-- (arquivos auxiliares gerados pelo grub-mkrescue)
\end{lstlisting}

\code{grub-mkrescue} cria uma imagem El Torito com o GRUB para BIOS
(\code{i386-pc}) e tamb\'em pode incluir suporte a EFI, dependendo dos
m\'odulos instalados. O projeto usa apenas o caminho BIOS.

Observa\c{c}\~oes da compila\c{c}\~ao de refer\^encia:

\begin{table}[H]
\centering
\caption{Artefatos gerados na compila\c{c}\~ao de refer\^encia.}
\label{tab:artefatos}
\begin{tabular}{@{}ll@{}}
\toprule
\textbf{Artefato} & \textbf{Valor observado} \\
\midrule
ELF & \code{target/x86\_64-unknown-none/release/vga-kernel} \\
Tamanho do ELF & aproximadamente 48 KiB \\
ISO & \code{vga-kernel.iso} \\
Tamanho da ISO & aproximadamente 33,8 MB \\
Setores da ISO & 16.517 setores de 2.048 bytes \\
Sistema de arquivos & ISO 9660 com Rock Ridge \\
Boot & El Torito, \code{boot\_hybrid.img} \\
\bottomrule
\end{tabular}
\end{table}

\section{Execu\c{c}\~ao}

\subsection{\code{make run}}

\begin{lstlisting}[style=shell,caption={Execu\c{c}\~ao padr\~ao.},label={lst:makerun}]
make run
# equivale a:
qemu-system-x86_64 -cdrom vga-kernel.iso -boot order=d -m 128M -serial stdio
\end{lstlisting}

A janela do QEMU mostra a tela VGA. O terminal de onde o comando foi
executado recebe a mesma sa\'ida pela serial e pode enviar comandos.

\subsection{Execu\c{c}\~ao sem interface gr\'afica}

Para testes automatizados, use \code{-display none} e sockets:

\begin{lstlisting}[style=shell,caption={QEMU headless com sockets de serial e monitor.},label={lst:qemuheadless}]
qemu-system-x86_64 -cdrom vga-kernel.iso -boot order=d -m 128M \
  -display none \
  -serial unix:/tmp/vga-serial.sock,server=on,wait=off \
  -monitor unix:/tmp/vga-mon.sock,server=on,wait=off
\end{lstlisting}

O monitor permite inspecionar a mem\'oria VGA:

\begin{lstlisting}[style=shell,caption={Dump da mem\'oria de texto VGA pelo monitor.},label={lst:vgadump}]
xp /4000bx 0xb8000
\end{lstlisting}

Cada c\'elula tem dois bytes: caractere e atributo. Um script pode reconstruir
as 25 linhas e verificar o conte\'udo.

\subsection{Valida\c{c}\~ao do ELF e da ISO}

\begin{lstlisting}[style=shell,caption={Verifica\c{c}\~oes \'uteis.},label={lst:validacoes}]
grub-file --is-x86-multiboot2 target/x86_64-unknown-none/release/vga-kernel
file vga-kernel.iso
xorriso -indev vga-kernel.iso -report_el_torito plain
readelf -h target/x86_64-unknown-none/release/vga-kernel
readelf -S target/x86_64-unknown-none/release/vga-kernel
\end{lstlisting}

\section{Solu\c{c}\~ao de problemas}

\begin{table}[H]
\centering
\caption{Problemas comuns e solu\c{c}\~oes.}
\label{tab:troubleshooting}
\begin{tabularx}{\linewidth}{@{}p{4.5cm}X@{}}
\toprule
\textbf{Sintoma} & \textbf{Causa prov\'avel / solu\c{c}\~ao} \\
\midrule
\code{can't find crate for `core'} & Cargo/rustc do sistema sem o alvo.
Use \code{rustup target add x86\_64-unknown-none} e
\code{rustup run stable cargo}. \\
\code{linker `ld.lld' not found} & Instale \code{lld} ou ajuste
\code{-C linker}. \\
\code{grub-mkrescue: command not found} & Instale \code{grub} e
\code{libisoburn}/\code{xorriso}. \\
QEMU tenta o disco e falha & Use \code{-boot order=d}. \\
Tela preta ap\'os o GRUB & Verifique se \code{gfxpayload=text} est\'a no
\code{grub.cfg} e se o ELF \'e Multiboot2 v\'alido. \\
Teclado n\~ao responde & Verifique \emph{legacy USB}/PS/2 na BIOS; teste a
serial. \\
Comandos aparecem na serial mas n\~ao no VGA & O modo texto VGA pode n\~ao
estar dispon\'ivel (UEFI sem CSM). \\
\code{make run} n\~ao abre janela & Falta servidor gr\'afico; use
\code{-display none} ou \code{-nographic}. \\
\bottomrule
\end{tabularx}
\end{table}

\section{Reprodutibilidade}

O build n\~ao depende de crates externos, mas depende de vers\~oes de Rust,
LLVM, GRUB e xorriso. Para maior reprodutibilidade, recomenda-se registrar:

\begin{itemize}
  \item a vers\~ao do Rust (\code{rustc --version});
  \item a vers\~ao do QEMU (\code{qemu-system-x86\_64 --version});
  \item a vers\~ao do GRUB (\code{grub-mkrescue --version});
  \item o hash do commit do projeto.
\end{itemize}

O \code{Cargo.lock} existe, mas n\~ao cont\'em depend\^encias externas; ele
apenas registra o pr\'oprio pacote.
